\ifdefined\gkver\else\def\gkver{student}\fi
\documentclass[\gkver]{gaokaozhenti}
\begin{document}

\chapter{2013年海南理综}
%% paperType: 理综物理部分

\begin{enumerate}
\item
%% number: 1
%% paperName: 2013年海南理综第 1 题 3 分
%% typeId: 1
%% type: 单选题
%% chapter: 电场
%% point: 电场强度
%% method: 无
%% score: 3
%% degree: 850
%% duplicateId: 0
%% body:
如图，电荷量为 $q_1$ 和 $q_2$ 的两个点电荷分别位于 P 点和 Q 点。已知在 P、Q 连线上某点 R 处的电场强度为零，且 $PR=2RQ$。则\xzanswer{B}

\onepicture[w=4.5cm, align=right]{figs/01.png}

\fourchoices[answer=B]
{$q_1=2q_2$}
{$q_1=4q_2$}
{$q_1=-2q_2$}
{$q_1=-4q_2$}

%% answer: B
%% memo:
\memoanswer{
由 $k\frac{q_1}{r_1^2}=k\frac{q_2}{r_2^2}$ 得：$q_1=4q_2$，故选 B。
}

\item
%% number: 2
%% paperName: 2013年海南理综第 2 题 3 分
%% typeId: 1
%% type: 单选题
%% chapter: 运动和力
%% point: 变加速直线运动
%% method: 无
%% score: 3
%% degree: 650
%% duplicateId: 0
%% body:
一质点受多个力的作用，处于静止状态，现使其中一个力的大小逐渐减小到零，再沿原方向逐渐恢复到原来的大小。在此过程中，其它力保持不变，则质点的加速度大小 $a$ 和速度大小 $v$ 的变化情况是\xzanswer{C}

\fourchoices[answer=C]
{$a$ 和 $v$ 都始终增大}
{$a$ 和 $v$ 都先增大后减小}
{$a$ 先增大后减小，$v$ 始终增大}
{$a$ 和 $v$ 都先减小后增大}

%% answer: C
%% memo:
\memoanswer{
由于质点所受合外力先增大后减小，故由牛顿第二定律知：质点做加速度先增大后减小的加速运动，所以 C 项正确。
}

\item
%% number: 3
%% paperName: 2013年海南理综第 3 题 3 分
%% typeId: 1
%% type: 单选题
%% chapter: 交流电
%% point: 交流电
%% method: 无
%% score: 3
%% degree: 650
%% duplicateId: 0
%% body:
通过一阻值 $R=100\UO$ 的电阻的交变电流如图所示，其周期为 $1\Us$。电阻两端电压的有效值为\xzanswer{B}

\onepicture[w=4.0cm, align=right]{figs/03.png}

\fourchoices[answer=B]
{$12\UV$}
{$4\sqrt{10}\UV$}
{$15\UV$}
{$8\sqrt{5}\UV$}

%% answer: B
%% memo:
\memoanswer{
设电流的有效值为 $I$，则有 $I^2RT=I_1^2R\left(\frac{4}{5}T\right)+I_2^2R\left(\frac{1}{5}T\right)$，其中 $I_1=0.1\UA$，$I_2=0.2\UA$，解得 $I=\frac{\sqrt{10}}{25}\UA$，电阻两端的电压 $U=IR=4\sqrt{10}\UV$，所以 B 项正确。
}

\item
%% number: 4
%% paperName: 2013年海南理综第 4 题 3 分
%% typeId: 1
%% type: 单选题
%% chapter: 直线运动
%% point: 运动图象
%% method: 无
%% score: 3
%% degree: 650
%% duplicateId: 0
%% body:
一物体做直线运动，其加速度随时间变化的 $a-t$ 图像如图所示。下列 $v-t$ 图像中，可能正确描述此物体运动的是\xzanswer{D}

\onepicture[w=3.5cm, align=right]{figs/04a.png}

\fourchoices[answer=D, ispicture=true, h=2.0cm]
{figs/04b.png}
{figs/04c.png}
{figs/04d.png}
{figs/04e.png}

%% answer: D
%% memo:
\memoanswer{
由 $a-t$ 图象可知：在 $0\sim\frac{T}{2}$ 内，物体做匀加速直线运动；$\frac{T}{2}\sim T$ 内，物体做匀速运动；$T\sim\frac{3T}{2}$ 内，物体做匀减速直线运动，且在 $\frac{3T}{2}$ 末速度为 0；$\frac{3T}{2}\sim 2T$ 内，物体做反向的匀加速运动；故 D 项正确。
}

\item
%% number: 5
%% paperName: 2013年海南理综第 5 题 3 分
%% typeId: 1
%% type: 单选题
%% chapter: 曲线运动
%% point: 人造地球卫星
%% method: 无
%% score: 3
%% degree: 650
%% duplicateId: 0
%% body:
“北斗”卫星导航定位系统由地球静止轨道卫星（同步卫星）、中轨道卫星和倾斜同步卫星组成。地球静止轨道卫星和中轨道卫星都在圆轨道上运行，它们距地面的高度分别约为地球半径的 6 倍和 3.4 倍，下列说法中正确的是\xzanswer{A}

\fourchoices[answer=A]
{静止轨道卫星的周期约为中轨道卫星的 2 倍}
{静止轨道卫星的线速度大小约为中轨道卫星的 2 倍}
{静止轨道卫星的角速度大小约为中轨道卫星的 1/7}
{静止轨道卫星的向心加速度大小约为中轨道卫星的 1/7}

%% answer: A
%% memo:
\memoanswer{
根据卫星绕地球做匀速圆周运动有 $G\frac{Mm}{r^2}=m\frac{v^2}{r}=m\omega^2r=m\left(\frac{2\pi}{T}\right)^2=ma_n$ 可知：$T_{\text{静}}=2T_{\text{中}}$，故 A 项正确。
}

\item
%% number: 6
%% paperName: 2013年海南理综第 6 题 3 分
%% typeId: 1
%% type: 单选题
%% chapter: 电磁感应
%% point: 楞次定律
%% method: 无
%% score: 3
%% degree: 850
%% duplicateId: 0
%% body:
如图，水平桌面上固定有一半径为 $R$ 的金属细圆环，环面水平，圆环每单位长度的电阻为 $r$，空间有一匀强磁场，磁感应强度大小为 $B$，方向竖直向下；一长度为 $2R$、电阻可忽略的导体棒置于圆环左侧并与环相切，切点为棒的中点。棒在拉力的作用下以恒定加速度 $a$ 从静止开始向右运动，运动过程中棒与圆环接触良好。下列说法正确的是\xzanswer{D}

\onepicture[w=3.5cm, align=right]{figs/06.png}

\fourchoices[answer=D]
{拉力的大小在运动过程中保持不变}
{棒通过整个圆环所用的时间为 $\sqrt{\frac{2R}{a}}$}
{棒经过环心时流过棒的电流为 $\frac{B\sqrt{2aR}}{\pi r}$}
{棒经过环心时所受安培力的大小为 $\frac{8B^2R\sqrt{2aR}}{\pi r}$}

%% answer: D
%% memo:
\memoanswer{
棒通过整个圆环所用的时间为 $t=\sqrt{\frac{4R}{a}}$，故 B 错；棒经过环心时棒上产生的感应电动势为 $E=2BR\sqrt{2aR}$，回路中的电阻为 $\frac{\pi Rr}{2}$，流过棒的电流为 $I=\frac{4B\sqrt{2aR}}{\pi r}$，此时棒所受安培力的大小为 $\frac{8B^2R\sqrt{2aR}}{\pi r}$，故 C 错 D 正确。
}

\item
%% number: 7
%% paperName: 2013年海南理综第 7 题 5 分
%% typeId: 2
%% type: 多选题
%% chapter: 其他
%% point: 物理学史
%% method: 无
%% score: 5
%% degree: 850
%% duplicateId: 0
%% body:
科学家关于物体运动的研究对树立正确的自然观具有重要作用。下列说法符合历史事实的是\xzanswer{BCD}

\fourchoices[answer=BCD]
{亚里士多德认为，必须有力作用在物体上，物体的运动状态才会改变}
{伽利略通过“理想实验”得出结论：运动必具有一定速度，如果它不受力，它将以这一速度永远运动下去}
{笛卡儿指出：如果运动中的物体没有受到力的作用，它将继续以同一速度沿同一直线运动，既不停下来也不偏离原来的方向}
{牛顿认为，物体具有保持原来匀速直线运动状态或静止状态的性质}

%% answer: BCD
%% memo:
\memoanswer{
本题原卷及材料中未提供详解，待补充。
}

\item
%% number: 8
%% paperName: 2013年海南理综第 8 题 5 分
%% typeId: 2
%% type: 多选题
%% chapter: 曲线运动
%% point: 圆周运动的应用
%% method: 无
%% score: 5
%% degree: 650
%% duplicateId: 0
%% body:
关于物体所受合外力的方向，下列说法正确的是\xzanswer{AD}

\fourchoices[answer=AD]
{物体做速率逐渐增加的直线运动时，其所受合外力的方向一定与速度方向相同}
{物体做变速率曲线运动时，其所受合外力的方向一定改变}
{物体做变速率圆周运动时，其所受合外力的方向一定指向圆心}
{物体做匀速率曲线运动时，其所受合外力的方向总是与速度方向垂直}

%% answer: AD
%% memo:
\memoanswer{
物体所受外力的方向与加速度的方向始终相同，与速度方向无关。当物体做加速直线运动时，合外力方向与速度方向相同，做减速直线运动时，合外力方向与速度方向相反；做匀变速曲线运动时，合外力方向不变，且与速度方向之间有夹角；做非匀变速曲线运动时，合外力方向时刻变化；做匀速圆周运动时，合外力方向始终指向圆心。故 AD 正确。
}

\item
%% number: 9
%% paperName: 2013年海南理综第 9 题 5 分
%% typeId: 2
%% type: 多选题
%% chapter: 磁场
%% point: 磁场的叠加
%% method: 无
%% score: 5
%% degree: 650
%% duplicateId: 0
%% body:
三条在同一平面（纸面）内的长直绝缘导线组成一等边三角形，在导线中通过的电流均为 $I$，方向如图所示。$a$、$b$ 和 $c$ 三点分别位于三角形的三个顶角的平分线上，且到相应顶点的距离相等。将 $a$、$b$ 和 $c$ 处的磁感应强度大小分别记为 $B_1$、$B_2$ 和 $B_3$，下列说法正确的是\xzanswer{AC}

\onepicture[w=3.8cm, align=right]{figs/09.png}

\fourchoices[answer=AC]
{$B_1=B_2<B_3$}
{$B_1=B_2=B_3$}
{$a$ 和 $b$ 处磁场方向垂直于纸面向外，$c$ 处磁场方向垂直于纸面向里}
{$a$ 处磁场方向垂直于纸面向外，$b$ 和 $c$ 处磁场方向垂直于纸面向里}

%% answer: AC
%% memo:
\memoanswer{
由磁场的叠加原理和安培定则知：$a$ 和 $b$ 处磁场方向垂直于纸面向外，$c$ 处磁场方向垂直于纸面向里，且 $B_1=B_2<B_3$。故 AC 正确。
}

\item
%% number: 10
%% paperName: 2013年海南理综第 10 题 5 分
%% typeId: 2
%% type: 多选题
%% chapter: 电磁感应
%% point: 楞次定律
%% method: 无
%% score: 5
%% degree: 650
%% duplicateId: 0
%% body:
如图，在水平光滑桌面上，两相同的矩形刚性小线圈分别叠放在固定的绝缘矩形金属框的左右两边上，且每个小线圈都各有一半面积在金属框内，在金属框接通逆时针方向电流的瞬间\xzanswer{BC}

\onepicture[w=6.0cm, align=right]{figs/10.png}

\fourchoices[answer=BC]
{两小线圈会有相互靠拢的趋势}
{两小线圈会有相互远离的趋势}
{两小线圈中感应电流都沿顺时针方向}
{左边小线圈中感应电流沿顺时针方向，右边小线圈中感应电流沿逆时针方向}

%% answer: BC
%% memo:
\memoanswer{
在金属框接通逆时针方向电流的瞬间，穿过两线圈的磁通量增加，为了阻碍磁通量的增加，由楞次定律知：左、右两线圈分别向左、右移动，两小线圈中感应电流都沿顺时针方向。故 BC 正确。
}

\item
%% number: 11
%% paperName: 2013年海南理综第 11 题 6 分
%% typeId: 4
%% type: 实验题
%% chapter: 直线运动
%% point: 打点计时器公式
%% method: 无
%% score: 6
%% degree: 650
%% duplicateId: 0
%% body:
某同学用图（a）所示的实验装置验证机械能守恒定律。已知打点计时器所用电源的频率为 $50\UHz$，当地重力加速度为 $g=9.80\Umsq$。实验中该同学得到的一条点迹清晰的完整纸带如图（b）所示。纸带上的第一个点记为 $O$，另选连续的三个点 $A$、$B$、$C$ 进行测量，图中给出了这三个点到 $O$ 点的距离 $h_A$、$h_B$ 和 $h_C$ 的值。回答下列问题（计算结果保留 3 位有效数字）

\twopicture[showlabel=false, wA=2.8cm, wB=5.5cm]
{figs/11a.png}
{figs/11b.png}

\begin{enumerate}
	\item
	打点计时器打 $B$ 点时，重物速度的大小 $v_B=$\tkanswer[1.5cm]{3.90}$\Ums$；

	\item
	通过分析该同学测量的实验数据，他的实验结果是否验证了机械能守恒定律？简要说明分析的依据。
\end{enumerate}

%% answer: 
\jdanswer{
\begin{enumerate}
	\item
	3.90

	\item
	设重物质量为 $m$，$OB$ 对应的下落过程中，重力势能减少量为 $mgh_B=7.70m\UJ$，动能增加量为 $\frac{1}{2}mv_B^2=7.61m\UJ$，在误差允许范围内，可以认为相等，因此验证了机械能守恒定律。

\end{enumerate}
}
%% memo:
\memoanswer{
（1）根据匀变速运动的规律知：$v_B=\frac{x_3-x_1}{2T}=3.90\Ums$。

（2）设重物质量为 $m$，$OB$ 对应的下落过程中，重力势能减少量为 $mgh_B=7.70m\UJ$，动能增加量为 $\frac{1}{2}mv_B^2=7.61m\UJ$，在误差允许范围内，可以认为相等，因此验证了机械能守恒定律。
}

\item
%% number: 12
%% paperName: 2013年海南理综第 12 题 9 分
%% typeId: 4
%% type: 实验题
%% chapter: 电路
%% point: 电路实验
%% method: 无
%% score: 9
%% degree: 650
%% duplicateId: 0
%% body:
某同学将量程为 $200\UmuA$、内阻为 $500\UO$ 的表头 $\mu$A 改装成量程为 $1\UmA$ 和 $10\UmA$ 的双量程电流表，设计电路如图（a）所示。定值电阻 $R_1=500\UO$，$R_2$ 和 $R_3$ 的值待定，$S$ 为单刀双掷开关，$A$、$B$ 为接线柱。回答下列问题：

\threepicture[showlabel=false, wA=3.8cm, wB=4.6cm, wC=4.6cm]
{figs/12a.png}
{figs/12b.png}
{figs/12c.png}

\begin{enumerate}
	\item
	按图（a）在图（b）中将实物连线；

	\item
	表笔 $a$ 的颜色为\tkanswer[1.5cm]{黑}色（填“红”或“黑”）

	\item
	将开关 $S$ 置于“1”挡时，量程为\tkanswer[1.5cm]{10}$\UmA$；

	\item
	定值电阻的阻值 $R_2=$\tkanswer[1.5cm]{25}$\UO$，$R_3=$\tkanswer[1.5cm]{225}$\UO$。（结果取 3 位有效数字）

	\item
	利用改装的电流表进行某次测量时，$S$ 置于“2”挡，表头指示如图（c）所示，则所测量电流的值为\tkanswer[1.5cm]{0.780}$\UmA$。
\end{enumerate}

%% answer: 
\jdanswer{
\begin{enumerate}
	\item
	如图（答图）所示

	\item
	黑

	\item
	10

	\item
	25，225

	\item
	0.780（0.78 同样给分）

\end{enumerate}
}
%% memo:
\memoanswer{
（1）按图（a）连接电路如答图所示。

\onepicture[w=6.0cm]{figs/12_an.jpg}

（2）表笔 $a$ 与表头“$-$”接线柱连接为黑色。

（3）开关 $S$ 置于“1”挡时，$R_3$ 为分流电阻小于开关 $S$ 置于“2”挡时的分流电阻 $R_3$、$R_2$ 之和，故为大量程 $10\UmA$。

（4）连接“1”时有：$(I_1-I_{\mu A})R_3=I_{\mu A}(R_{\mu A}+R_1+R_2)$，其中 $I_{\mu A}=200\UmuA$，$I_1=10\UmA$；连接“2”时有：$(I_2-I_{\mu A})(R_3+R_2)=I_{\mu A}(R_{\mu A}+R_1)$，$I_2=1\UmA$，联立代入数据解得 $R_2=25\UO$，$R_3=225\UO$。

（5）将表头刻度换成 $0\sim1\UmA$ 刻度值，可读的 $0.780\UmA$。
}

\item
%% number: 13
%% paperName: 2013年海南理综第 13 题 10 分
%% typeId: 6
%% type: 计算题
%% chapter: 功和能
%% point: 动能定理
%% method: 无
%% score: 10
%% degree: 650
%% duplicateId: 0
%% body:
一质量 $m=0.6\Ukg$ 的物体以 $v_0=20\Ums$ 的初速度从倾角为 $30^\circ$ 的斜坡底端沿斜坡向上运动。当物体向上滑到某一位置时，其动能减少了 $\Delta E_k=18\UJ$，机械能减少了 $\Delta E=3\UJ$。不计空气阻力，重力加速度 $g=10\Umsq$，求：

\begin{enumerate}
	\item
	物体向上运动时加速度的大小；

	\item
	物体返回斜坡底端时的动能。
\end{enumerate}

%% answer: 
\jdanswer{
\begin{enumerate}
	\item
	$6\Umsq$

	\item
	$80\UJ$

\end{enumerate}
}
%% memo:
\memoanswer{
（1）设物体在运动过程中所受的摩擦力大小为 $f$，向上运动的加速度大小为 $a$，由牛顿第二定律有

$a=\frac{mg\sin\alpha+f}{m}$ ①

设物体动能减少 $\Delta E_k$ 时，在斜坡上运动的距离为 $s$，由动能定理可得

$\Delta E_k=(mg\sin\alpha+f)s$ ②

$\Delta E=fs$ ③

联立①②③式并带入数据可得：

$a=6\Umsq$ ④

（2）设物体沿斜坡向上运动的最大距离为 $s_m$，由运动学规律可得

$s_m=\frac{v_0^2}{2a}$ ⑤

设物体返回底端时的动能为 $E_k$，由动能定理有

$E_k=(mg\sin\alpha-f)s_m$ ⑥

联立①④⑤⑥式并代入数据可得

$E_k=80\UJ$ ⑦
}

\item
%% number: 14
%% paperName: 2013年海南理综第 14 题 13 分
%% typeId: 6
%% type: 计算题
%% chapter: 磁场
%% point: 带电粒子在磁场中的运动
%% method: 无
%% score: 13
%% degree: 650
%% duplicateId: 0
%% body:
如图，纸面内有 $E$、$F$、$G$ 三点，$\angle GEF=30^\circ$，$\angle EFG=135^\circ$。空间有一匀强磁场，磁感应强度大小为 $B$，方向垂直于纸面向外。先使带有电荷量为 $q$（$q>0$）的点电荷 $a$ 在纸面内垂直于 $EF$ 从 $F$ 点射出，其轨迹经过 $G$ 点；再使带有同样电荷量的点电荷 $b$ 在纸面内与 $EF$ 成一定角度从 $E$ 点射出，其轨迹也经过 $G$ 点。两点电荷从射出到经过 $G$ 点所用的时间相同，且经过 $G$ 点时的速度方向也相同。已知点电荷 $a$ 的质量为 $m$，轨道半径为 $R$，不计重力。求：

\begin{enumerate}
	\item
	点电荷 $a$ 从射出到经过 $G$ 点所用的时间；

	\item
	点电荷 $b$ 的速度大小。
\end{enumerate}

\onepicture[w=4.5cm, align=right]{figs/14.png}

%% answer: 
\jdanswer{
\begin{enumerate}
	\item
	$\frac{\pi m}{2qB}$

	\item
	$\frac{4qBR}{3m}$

\end{enumerate}
}
%% memo:
\memoanswer{
（1）设点电荷 $a$ 的速度大小为 $v$，由牛顿第二定律得

$qvB=m\frac{v^2}{R}$ ①

由①式得：

$v=\frac{qBR}{m}$ ②

设点电荷 $a$ 做圆周运动的周期为 $T$，有

$T=\frac{2\pi m}{qB}$ ③

如图，$O$ 和 $O_1$ 分别是 $a$ 和 $b$ 的圆轨道的圆心。设 $a$ 在磁场中偏转的角度为 $\theta$，由几何关系得：

$\theta=90^\circ$ ④

故 $a$ 从开始运动到经过 $G$ 点所用时间 $t$ 为

$t=\frac{\pi m}{2qB}$ ⑤

（2）设点电荷 $b$ 的速度大小为 $v_1$，轨道半径为 $R_1$，$b$ 在磁场中偏转的角度为 $\theta_1$，依题意有

$t=\frac{R_1\theta_1}{v_1}=\frac{R\theta}{v}$ ⑥

由⑥式得

$v_1=\frac{R_1\theta_1}{R\theta}v$ ⑦

由于两轨道在 $G$ 点相切，所有过 $G$ 点的半径 $OG$ 和 $O_1G$ 在同一直线上，由几何关系和题给条件得

$\theta_1=60^\circ$ ⑧

$R_1=2R$ ⑨

联立②④⑦⑧⑨式，解得

$v_1=\frac{4qBR}{3m}$ ⑩
}

\item
%% number: 15
%% paperName: 2013年海南理综第 15 题 8 分
%% typeId: 6
%% type: 计算题
%% chapter: 气体性质
%% point: 玻意耳定律
%% method: 无
%% score: 8
%% degree: 650
%% duplicateId: 0
%% body:
[选修 3-3]（12 分）

\begin{enumerate}
	\item
	（4 分）下列说法正确的是\xzanswer{ACD}

	\fivechoices[answer=ACD]
	{把一枚针轻放在水面上，它会浮在水面，这是由于水表面存在表面张力的缘故}
	{水在涂有油脂的玻璃板上能形成水珠，而在干净的玻璃板上却不能，为是因为油脂使水的表面张力增大的缘故}
	{在围绕地球飞行的宇宙飞船中，自由飘浮的水滴呈球形，这是表面张力作用的结果}
	{在毛细现象中，毛细管中的液面有的升高，有的降低，这与液体的种类和毛细管的材质有关}
	{当两薄玻璃板间夹有一层水膜时，在垂直于玻璃板的方向很难将玻璃板拉开，这是由于水膜具有表面张力的缘故}

	\item
	（8 分）如图，一带有活塞的气缸通过底部的水平细管与一个上端开口的竖直管相连，气缸与竖直管的横截面面积之比为 $3:1$，初始时，该装置的底部盛有水银；活塞与水银面之间有一定量的气体，气柱高度为 $l$（以 $\Ucm$ 为单位）；竖直管内的水银面比气缸内的水银面高出 $\frac{3}{8}l$。现使活塞缓慢向上移动 $\frac{11}{32}l$，这时气缸和竖直管内的水银面位于同一水平面上，求初始时气缸内气体的压强（以 $\UcmHg$ 为单位）

	\onepicture[w=3.2cm, align=right]{figs/15.png}
\end{enumerate}

%% answer: 
\jdanswer{
\begin{enumerate}
	\item
	ACD

	\item
	$\frac{15}{8}l$

\end{enumerate}
}
%% memo:
\memoanswer{
（2）设 $S$ 为气缸的横截面积，$P$ 为活塞处于初始位置时气缸内气体的压强，$P_0$ 为大气压强，有

$P=P_0+\frac{3}{8}l$ ①

在活塞上移 $\frac{11}{32}l$ 后，气缸内气体的压强变为 $P_0$，设气体的体积为 $V'$，由玻意耳定律得

$P_0V'=PSl$ ②

设气缸内水银面上升 $\Delta x$，有

$\Delta x=\frac{3}{8}l-3\Delta x$ ③

$V'=\left(l+\frac{11}{32}l-\Delta x\right)S$ ④

联立①②③④式，解得

$P=\frac{15}{8}l$ ⑤
}

\item
%% number: 16
%% paperName: 2013年海南理综第 16 题 8 分
%% typeId: 6
%% type: 计算题
%% chapter: 光学
%% point: 光的折射
%% method: 无
%% score: 8
%% degree: 650
%% duplicateId: 0
%% body:
[选修 3-4]（12 分）

\begin{enumerate}
	\item
	（4 分）下列选项与多普勒效应有关的是\xzanswer{BDE}

	\fivechoices[answer=BDE]
	{科学家用激光测量月球与地球间的距离}
	{医生利用超声波探测病人血管中血液的流速}
	{技术人员用超声波探测金属、陶瓷、混凝土中是否有气泡}
	{交通警察向车辆发射超声波并通过测量反射波的频率确定车辆行进的速度}
	{科学家通过比较星球与地球上同种元素发出光的频率来计算星球远离地球的速度}

	\item
	（8 分）如图，三棱镜的横截面为直角三角形 $ABC$，$\angle A=30^\circ$，$AC$ 平行于光屏 $MN$，与光屏的距离为 $L$。棱镜对红光的折射率为 $n_1$，对紫光的折射率为 $n_2$。一束很细的白光由棱镜的侧面 $AB$ 垂直射入，直接到达 $AC$ 面并射出。画出光路示意图，并标出红光和紫光射在光屏上的位置，求红光和紫光在光屏上的位置之间的距离。

	\onepicture[w=3.5cm, align=right]{figs/16.png}
\end{enumerate}

%% answer: 
\jdanswer{
\begin{enumerate}
	\item
	BDE

	\item
	$L\left(\frac{n_2}{\sqrt{4-n_2^2}}-\frac{n_1}{\sqrt{4-n_1^2}}\right)$

\end{enumerate}
}
%% memo:
\memoanswer{
（2）光路图所示，红光和紫光在 $AC$ 面上的入射角相同，设为 $i$，折射角分别为 $r_1$ 和 $r_2$，它们射到屏上的位置离 $O$ 点的距离分别为 $d_1$ 和 $d_2$。由折射定律得

$n_1\sin i=\sin r_1$ ①

$n_2\sin i=\sin r_2$ ②

由几何关系得

$i=\angle A$ ③

$d_1=L\tan r_1$ ④

$d_2=L\tan r_2$ ⑤

联立①②③④⑤式并利用题给条件得，红光和紫光在光屏上的位置之间的距离为

$d_2-d_1=L\left(\frac{n_2}{\sqrt{4-n_2^2}}-\frac{n_1}{\sqrt{4-n_1^2}}\right)$ ⑥

\onepicture[w=5.0cm]{figs/16_an.jpg}
}

\item
%% number: 17
%% paperName: 2013年海南理综第 17 题 8 分
%% typeId: 6
%% type: 计算题
%% chapter: 运动和力
%% point: 动量守恒定律
%% method: 无
%% score: 8
%% degree: 650
%% duplicateId: 0
%% body:
[选修 3-5]（12 分）

\begin{enumerate}
	\item
	（4 分）原子核 \ce{^{232}_{90}Th} 具有天然放射性，它经过若干次 $\alpha$ 衰变和 $\beta$ 衰变后会变成新的原子核。下列原子核中，有三种是 \ce{^{232}_{90}Th} 衰变过程中可以产生的，它们是\xzanswer{ACD}

	\fivechoices[answer=ACD]
	{\ce{^{204}_{82}Pb}}
	{\ce{^{203}_{82}Pb}}
	{\ce{^{216}_{84}Po}}
	{\ce{^{224}_{88}Ra}}
	{\ce{^{226}_{88}Ra}}

	\item
	（8 分）如图，光滑水平面上有三个物块 $A$、$B$ 和 $C$，它们具有相同的质量，且位于同一直线上。开始时，三个物块均静止，先让 $A$ 以一定速度与 $B$ 碰撞，碰后它们粘在一起，然后又一起与 $C$ 碰撞并粘在一起，求前后两次碰撞中损失的动能之比。

	\onepicture[w=5.0cm, align=right]{figs/17.png}
\end{enumerate}

%% answer: 
\jdanswer{
\begin{enumerate}
	\item
	ACD

	\item
	3

\end{enumerate}
}
%% memo:
\memoanswer{
（2）设三个物块 $A$、$B$ 和 $C$ 的质量均为 $m$，$A$ 与 $B$ 碰撞前 $A$ 的速度为 $v$，碰撞后共同速度为 $v_1$，$A$、$B$ 与 $C$ 碰撞后共同速度为 $v_2$，由动量守恒定律得

$mv=(m+m)v_1$ ①

$mv=(m+m+m)v_2$ ②

设第一次碰撞中动能的损失为 $\Delta E_1$，第二次碰撞中动能的损失为 $\Delta E_2$，由能量守恒定律得

$\frac{1}{2}mv^2=\frac{1}{2}(m+m)v_1^2+\Delta E_1$ ③

$\frac{1}{2}(m+m)v_1^2=\frac{1}{2}(m+m+m)v_2^2+\Delta E_2$ ④

联立①②③④解得

$\frac{\Delta E_1}{\Delta E_2}=3$ ⑤
}

\end{enumerate}

\end{document}
